% Documentary recovery for Article V; RESULTADO_RECUPERADO.
% Source owner: /Users/ruben/Documents/New project/output/REVISION_EDITORIAL_HMT_MD_20260905/01_FUENTES/INTEGRAL/tree/New project/output/ENSAMBLADO_CAUSAL_RESTAURADO_HMT_MD_20260905/fuente/manuscrito/sections/md/delta_127/04b_completaciones_cuanticas.tex
% Source lines: 1--334.
% REV02: algebraic construction; the statistical application is consolidated in 40.
% Editorial correspondence: gestion/CONSOLIDACION_PAULI_FOCK_REV02.json.
\section{Algebraic construction of quantum completions over APP}
\label{v:rec:chap:completaciones-cuanticas}
\label{v:sec:completaciones-algebraicas}

APP and TPK provide oriented states, transport, and sheet parity.
Their linear realization enables construction of the symmetric and exterior algebras,
with their creation, annihilation, and occupation operators. This section
establishes these algebraic constructions, the cellular field sector, and
reversible measurement coupling. Application of the occupation spectra to equilibrium is developed in Section~\ref{v:rec:ocupacion:section:02}.
Each stage retains the conditions of the representation it uses.

\subsection{Linear state space associated with a particle realization}
\label{v:rec:completaciones:section:02}

Let \(S\) be a finite set of TPK route classes at the chosen depth,
and let
\[
 \mathcal H_1=\ell^2(S)
\]
be the Hilbert space with orthonormal basis \((e_s)_{s\in S}\). A choice of
positive weights \(E_s\) defines the diagonal operator
\[
 He_s=E_se_s.
\]
The choice of \(S\), its quotient by symmetries, and the weights are part
of the model. Once these are fixed, the multiparticle completions are canonical.

\subsection{Symmetric and exterior completions}
\label{v:rec:completaciones:section:03}

Define
\[
 \mathcal F_+(\mathcal H_1)=
 \bigoplus_{n\geq0}\operatorname{Sym}^n\mathcal H_1,
 \qquad
 \mathcal F_-(\mathcal H_1)=
 \bigoplus_{n\geq0}\bigwedge^n\mathcal H_1.
\]
The first is the bosonic completion; the second is the fermionic completion. In both,
the degree-zero summand is the vacuum space. The exterior powers
are equipped with the inner product for which ordered wedges of vectors
from an orthonormal basis are orthonormal; the symmetric powers are used
with the usual normalization of creation and annihilation operators.

\subsection{Sheet grading and exchange algebra}
\label{v:rec:completaciones:section:04}

Suppose that the oriented sheet defines an additive character
\[
 p:\operatorname{Mor}(\mathsf P_{\rm TPK})\longrightarrow\mathbb Z/2\mathbb Z,
 \qquad p(\gamma\eta)=p(\gamma)+p(\eta).
\]
On the corresponding graded linear space, introduce the symmetry
\[
 \tau(v\otimes w)=(-1)^{p(v)p(w)}w\otimes v.
\]

\begin{theorem}[Statistical completion determined by parity]
\label{v:rec:completaciones:statement:02}
The quotient of the tensor algebra by the relations
\[
 v\otimes w-(-1)^{p(v)p(w)}w\otimes v
\]
is symmetric on the even sector and exterior on the odd sector. Two identical
odd states have zero product; even states admit arbitrary
occupation.
\end{theorem}

\begin{proof}
If \(p(v)=p(w)=0\), the relation is \(vw=wv\). If
\(p(v)=p(w)=1\), it is \(vw=-wv\); setting \(v=w\) and working in
characteristic zero gives \(v^2=0\). Mixed sectors commute without
an additional sign. This is the universal property of the graded symmetric
algebra.
\end{proof}

\subsection{Canonical operators and idempotence of occupation}
\label{v:pauli:car-construccion}

The parity character determines the completion; the product of that
completion then determines the operator algebra. In the
exterior sector, creation \(a_s^\dagger\) is exterior multiplication by
\(e_s\), and \(a_s\) is its adjoint.

\begin{theorem}[Pauli exclusion principle in the exterior completion]
\label{v:rec:completaciones:statement:01}
\label{v:rec:thm:principio-exclusion-pauli}
For every mode \(s\in S\), the creation and annihilation operators of the
exterior completion satisfy
\begin{equation}
 \{a_s,a_t^\dagger\}=\delta_{st}I,
 \qquad
 \{a_s,a_t\}=\{a_s^\dagger,a_t^\dagger\}=0.
 \label{v:pauli:eq:car}
\end{equation}
In particular, the occupation operator
\(N_s=a_s^\dagger a_s\) is a projector:
\begin{equation}
 N_s^2=N_s,
 \qquad \Spec(N_s)\subseteq\{0,1\}.
 \label{v:pauli:eq:proyeccion-ocupacion}
\end{equation}
The eigenvalues \(0\) and \(1\) correspond, respectively, to the absence
and presence of mode \(s\).
\end{theorem}

\begin{proof}
On an ordered wedge, \(a_s\) removes the factor \(e_s\), when present,
with the sign determined by its position; when absent, it gives zero. Composing
this operation with exterior multiplication proves relations
\eqref{v:pauli:eq:car}: for \(s=t\), the two compositions cover the
complementary cases of absence and presence; for \(s\ne t\), the exchange
signs cancel. Antisymmetry implies
\(e_s\wedge e_s=0\) and
\(a_s^2=(a_s^\dagger)^2=0\). Therefore
\begin{equation}
 N_s^2=a_s^\dagger a_s a_s^\dagger a_s
 =a_s^\dagger(I-a_s^\dagger a_s)a_s
 =a_s^\dagger a_s=N_s.
 \label{v:pauli:eq:prueba-idempotencia}
\end{equation}
The operator is self-adjoint, and every eigenvalue of an idempotent satisfies
\(\lambda^2=\lambda\). Hence its spectrum is contained in
\(\{0,1\}\). The vacuum and the state \(e_s\) realize both values
\cite{Pauli1925,Dirac1926}.
\end{proof}

In the symmetric sector, the bosonic operators \(b_s^\dagger,b_s\)
act on the normalized occupation basis with factors
\(\sqrt{m_s+1}\) and \(\sqrt{m_s}\), respectively. Their compositions
give
\begin{equation}
 [b_s,b_t^\dagger]=\delta_{st}I,
 \qquad [b_s,b_t]=[b_s^\dagger,b_t^\dagger]=0.
 \label{v:pauli:eq:ccr}
\end{equation}
These identities are understood on the algebraic finite-particle domain,
which is invariant under the operators considered.

Theorem~\ref{v:rec:completaciones:statement:02} turns a compositional TPK parity into an exact
occupation rule. To obtain the physical spin--statistics theorem, one must
also construct a positive local Lorentzian theory and prove that
sheet parity agrees with spin under rotations. This equality does not
follow merely from both objects taking values modulo two.
